%!TEX root = ../../../../exa-ma-d7.1.tex

\section{Weak scaling of an adaptive Burgers solver}
\label{sec:app:specs:app-samurai-burgers-2d}

This mini-app measures the weak scalability of the distributed adaptive mesh of
samurai (see \Cref{sec:Samurai:software}) on the two-dimensional inviscid
Burgers equation. The test case is built so that every MPI process carries
exactly the same adapted mesh and the same amount of work. Any loss of
efficiency as the number of processes grows therefore comes from the
communications and from the distributed mesh adaptation, not from load
imbalance.

\Cref{tab:app-samurai-burgers-2d} describes the specifications of the
application.

\begin{table}[ht]
    \centering
    \begin{tblr}{
        colspec = {l X[10cm]},
        row{odd} = {numpexlightergray},
        hlines = {0.1pt, numpexgray},
        vlines = {numpexgray},
        row{1} = {numpexgray, fg=white, font=\bfseries},
    }
        Field & Details \\
        id & \texttt{app-samurai-burgers-2d} \\
        name & Adaptive Burgers 2D, weak scaling \\
        Partners & IPP \\
        PC & PC1 - ExaMA \\
        Responsible (Permanent) & L. Gouarin \\
        WP7 Engineer & \todo{To be confirmed} \\
        work\_package & WP1 \\
        application\_type & mini-app \\
        purpose & Weak scalability of the distributed multiresolution mesh adaptation on a load-balanced nonlinear transport problem \\
        Method-Algorithm WP1 & adaptive Cartesian mesh (interval representation), multiresolution analysis, finite volume, OSMP scheme, Strang splitting \\
        Method-Algorithm WP2 & \\
        Method-Algorithm WP3 & \\
        Method-Algorithm WP4 & \\
        Method-Algorithm WP5 & \\
        Method-Algorithm WP6 & \\
        WP7 & \\
        metrics & \SetCell{l} \texttt{weak-scalability}, \texttt{time-to-solution}, \texttt{regression-test-pass} \\
        status & benchmark-in-progress (repository-only, not yet in the operational Sheet) \\
        Benchmark scope & multi-node, full-system \\
        Framework & samurai \\
        parallel\_framework & MPI \\
        repo\_url & \url{https://github.com/hpc-maths/samurai}\\
    \end{tblr}
    \caption{Description of the mini-app \texttt{app-samurai-burgers-2d}.}
    \label{tab:app-samurai-burgers-2d}
\end{table}


%%%%%%%%%%%%%%%%%%%%%%%%%%%%%%%%%%%%%%%%%%%%%%%%%%%%%%%%%%%%%%%%%%%%%%%%%%%%%%%%%%%%%%%%%%%%%%%%%%%%%%%%%%%%%%%%%%%%%%%%


\subsection{Description of the benchmark}

We solve the two-dimensional inviscid Burgers equation
\begin{equation}\label{eq:specs:samurai:burgers}
  \partial_t u + \partial_x \left(\frac{u^2}{2}\right) + \partial_y \left(\frac{u^2}{2}\right) = 0,
\end{equation}
with periodic boundary conditions. On the reference box $[-1,1]^2$, the
initial condition is a Gaussian bump
\begin{equation}\label{eq:specs:samurai:burgers:init}
  u_0(x,y) = \frac{1}{2} \exp\left(-50\left[(x-x_c)^2 + (y-y_c)^2\right]\right),
\end{equation}
whose centre $(x_c, y_c)$ is shifted by $0.25$ from the centre of the box in
each direction. The nonlinear flux steepens the bump into a shock, so the
adapted mesh follows a moving, localized discontinuity until the final time
$T_f$.

\paragraph{Construction of the weak scaling.}
With $P = n_x \times n_y$ MPI processes, the global domain is a grid of
$n_x \times n_y$ copies of the reference box. The process of rank $r$ owns the
copy translated by $2\,(r \bmod n_x, \lfloor r / n_x \rfloor)$, and its
initial condition is \eqref{eq:specs:samurai:burgers:init} expressed relative
to its own box. Since the global problem is periodic with the period of the
reference box, every subdomain holds the same solution and the same adapted
mesh, up to a translation. The work per process is therefore constant by
construction, and the ideal weak scaling is a constant time to solution.

At every time step, the program compares each subdomain with its MPI neighbours,
level by level, and stops with an error if the adapted meshes differ. This
check verifies that the parallel mesh adaptation gives the same result as
on a single process, and that the load stays balanced throughout the run.

\paragraph{Numerical method.}
The mesh is adapted at every time step by multiresolution analysis
\cite{cohen_fully_2003,bellotti_multidimensional_2022}, with threshold
$\varepsilon = 10^{-3}$, between the levels $\ell_{\min}$ and $\ell_{\max}$
(4 and 10 by default), with a graduation width of 2. The flux is discretized
with the one-step monotonicity-preserving (OSMP) scheme of Daru and Tenaud
\todo{Add Daru \& Tenaud, J. Comput. Phys. (2004) to the Zotero library and cite it here.}
on a four-cell stencil. The two directions are coupled by a Strang splitting,
$X(\Delta t/2)\,Y(\Delta t)\,X(\Delta t/2)$. The time step is fixed for the
whole run from the finest level the mesh may reach,
$\Delta t = \mathrm{CFL}\,\Delta x_{\ell_{\max}} / \max|u_0|$ with
$\mathrm{CFL} = 0.95$, so that it stays the same whatever the number of
processes.


%%%%%%%%%%%%%%%%%%%%%%%%%%%%%%%%%%%%%%%%%%%%%%%%%%%%%%%%%%%%%%%%%%%%%%%%%%%%%%%%%%%%%%%%%%%%%%%%%%%%%%%%%%%%%%%%%%%%%%%%


\subsection{Benchmarking tools used}

The benchmark is the demo \texttt{demos/FiniteVolume/burgers\_os\_2d\_mpi.cpp}
of samurai, built with \texttt{SAMURAI\_WITH\_MPI=ON}. The process grid is set
with \texttt{-{}-npx} and \texttt{-{}-npy}, and the program refuses to run if
$n_x n_y \neq P$. A run on 8 processes reads
\begin{verbatim}
mpiexec -n 8 ./finite-volume-burgers-os-2d-mpi --npx=4 --npy=2 \
        --min-level=4 --max-level=10 --Tf=0.5 --no-output --timers
\end{verbatim}

The \texttt{-{}-timers} option prints the built-in timers of samurai at the end
of the run. They split the wall-clock time between mesh adaptation, ghost
updates and the numerical scheme. \texttt{-{}-no-output} removes file writing
from the measurement.

Let $T(P)$ be the time to solution on $P$ processes, taken on the slowest rank.
We report the weak scaling efficiency
\begin{equation}\label{eq:specs:samurai:burgers:weak-efficiency}
  E_w(P) = \frac{T(1)}{T(P)},
\end{equation}
for the whole time loop and for each of the timed components. Since the work
per process does not depend on $P$, the difference between $E_w(P)$ and 1
measures the parallel overhead directly.

The same program runs in the continuous integration of samurai on 2, 4, 8 and 9
processes, with several process-grid shapes ($2\times1$, $1\times2$, $4\times2$,
$2\times4$, $3\times3$). This gives a \texttt{regression-test-pass} metric for
the distributed mesh adaptation at small scale.


%%%%%%%%%%%%%%%%%%%%%%%%%%%%%%%%%%%%%%%%%%%%%%%%%%%%%%%%%%%%%%%%%%%%%%%%%%%%%%%%%%%%%%%%%%%%%%%%%%%%%%%%%%%%%%%%%%%%%%%%


\subsection{Input/Output Dataset Description}


\subsubsection{Input Data:}
  \begin{itemize}
  \item No external dataset: the domain, the initial condition and the mesh are
    generated by the program from its command-line parameters.
  \item Parameters: process grid (\texttt{-{}-npx}, \texttt{-{}-npy}), mesh levels
    (\texttt{-{}-min-level}, \texttt{-{}-max-level}), final time (\texttt{-{}-Tf}),
    CFL number (\texttt{-{}-cfl}) and multiresolution threshold
    (\texttt{-{}-mr-eps}).
  \item Software: samurai commit \texttt{1894c71}. The commit of the runs will
    be pinned with the results.
  \end{itemize}

\subsubsection{Output Data:}
  \begin{itemize}
  \item The timers report of each run, with the number of processes and the
    process grid.
  \item Optionally, the solution and the level of each cell in HDF5/XDMF
    format (\texttt{-{}-nfiles}), to check the adapted mesh visually.
  \end{itemize}


%%%%%%%%%%%%%%%%%%%%%%%%%%%%%%%%%%%%%%%%%%%%%%%%%%%%%%%%%%%%%%%%%%%%%%%%%%%%%%%%%%%%%%%%%%%%%%%%%%%%%%%%%%%%%%%%%%%%%%%%


\subsection{Results summary}

No scaling result is reported in this version of the deliverable. The next
version will give, for each run, the machine (one of those listed in Appendix A),
the compiler and MPI versions, the number of processes per node, the mesh
levels, the final time and the number of repetitions. It will show the time to
solution of each timed component and the efficiency $E_w(P)$ of
\eqref{eq:specs:samurai:burgers:weak-efficiency} against $P$.


%%%%%%%%%%%%%%%%%%%%%%%%%%%%%%%%%%%%%%%%%%%%%%%%%%%%%%%%%%%%%%%%%%%%%%%%%%%%%%%%%%%%%%%%%%%%%%%%%%%%%%%%%%%%%%%%%%%%%%%%


\subsection{Challenges identified}

\begin{itemize}
\item The mesh comparison between neighbours runs inside the time loop and ends
  with a global reduction. Its cost has to be timed separately, or the check
  disabled for the timing runs, so that it is not counted as overhead of the
  mesh adaptation.
\item The test case balances the load by construction. It measures the cost of
  the communications and of the distributed adaptation, but not the effect of
  an unbalanced adapted mesh. The double Mach reflection benchmark
  (\Cref{sec:app:specs:app-samurai-euler-dmr}) covers that case.
\end{itemize}
